\documentclass[10pt,twocolumn]{article}
\usepackage[letterpaper,hmargin=0.5in,vmargin=0.55in]{geometry}
\setlength{\columnsep}{0.25in}              % gutter between columns
\linespread{1.08}                           % a little more space between lines
\raggedbottom                               % no stretched gaps at column bottoms
\usepackage[T1]{fontenc}
\usepackage[dvipsnames,table]{xcolor}
\usepackage[all]{xy}
\usepackage[authoryear]{natbib}
\setlength{\bibhang}{1.5em}   % hanging indent size
\setlength{\bibsep}{4pt}      % space between entries
\usepackage{amsmath,amsthm,amssymb,latexsym}
\usepackage{graphicx}
\usepackage{booktabs}
\usepackage{float}
\usepackage{tabularx}                      % tables that fill the text width
\newcolumntype{L}{>{\raggedright\arraybackslash}X}
\usepackage{enumitem}
\setlist{nosep, leftmargin=*}              % compact bullet and numbered lists
\setlist[description]{nosep, leftmargin=1.2em}
\usepackage[font=small, skip=4pt]{caption} % smaller captions, closer to tables
\setlength{\intextsep}{8pt}                % space around [H] tables and figures
\setlength{\floatsep}{8pt}
\setlength{\textfloatsep}{10pt}
\AtBeginDocument{%
  \setlength{\abovedisplayskip}{6pt}\setlength{\belowdisplayskip}{6pt}%
  \setlength{\abovedisplayshortskip}{4pt}\setlength{\belowdisplayshortskip}{4pt}}
\usepackage{comment}
\usepackage{tikz}
\usepackage{tikz-cd}
\usetikzlibrary{automata, positioning, arrows}
\usepackage{microtype}                     % better line breaking in narrow columns
\setlength{\emergencystretch}{1.5em}
\usepackage{xurl}                          % let long URLs break anywhere
\usepackage{hyperref}

\begin{document}

% Title and abstract span both columns
\twocolumn[{
\begin{center}
    {\LARGE \textbf{Which Counties Will Flood Damage Hit, and How Much Will It Cost?\\[0.15cm]
    Forecasting County-Month Property Losses}} \\[0.3cm]
    {\normalsize Tatyana Sacdalan, Enci Zheng} \\[0.1cm]
    {\normalsize November 23, 2026}
\end{center}

\hrulefill

\begin{center}
    {\large \textbf{Abstract}}
\end{center}

{\small
Emergency responders need to know in advance which communities a flood will hit hardest, but
damage depends on the community as much as the storm: how much is built there, how close it
sits to rivers, and how well prepared it is. We predict monthly flood property damage for the
1,055 counties of 12 Midwest states from 2000 to 2023, combining SHELDUS loss records with rain,
snow, soil, land-cover, terrain and built-up data. A two-part hurdle model predicts whether any
damage occurs and, if it does, which loss class it falls in: under \$10k, \$10k--\$100k,
\$100k--\$1M, or \$1M or more. LightGBM, the main model, ranks county-months well on the test
years (ROC-AUC 0.89), ahead of a generalized linear model and each county's own damage history.
A separate causal model shows that 25\,mm more rain roughly doubles a county's chance of flood damage, that this effect holds after accounting for climate and preparedness, and that higher-income counties are damaged less often, while better-insured ones are not.
\par}

\vspace{0.4cm}
}]

\section{Introduction}

Floods cause large property losses across the Midwest, while emergency resources are limited.
To set priorities across counties, it helps to forecast the impact of a flood on each county,
not only how much rain will fall. We build an impact model from past floods and use it to rank
counties by their predicted chance of a severe loss.

Two features of the problem make this difficult. First, loss records are sparse and uneven:
only 3\% of county-months (9,845 of 303,840) have a flood loss in SHELDUS, reporting rates differ between states,
and losses spanning several counties are often split equally among them. Second, the effect of
storms on loss may be confounded. Wetter counties receive larger storms but may also adapt more,
through flood insurance and floodplain management, which lowers the damage a storm causes.

We address the first with a two-part model that separates whether damage occurs from how large
it is, trained on 2000--2016 and tested on 2017--2023. We address the second with a set of
causal models that test whether accounting for rain climate and measured adaptation changes
the estimated effect of storms on damage. Our main contribution is a county-month model that ranks counties by their chance of a severe flood loss; a secondary analysis tests whether climate and preparedness change the estimated storm effect, and estimates how much more rain raises a county's chance of flood damage.


\section{Related Works}

\paragraph{Floods, rainfall and damage.}
Heavy rainfall alone does not explain flooding. Across 420 U.S. catchments,
\citet{berghuijs2016} found that extreme precipitation controls floods only in the central arid region; elsewhere, rain on already wet soil or snowmelt dominates. \citet{slater2016} found rising flood frequency in the upper Midwest and Great Lakes, tracking long-term basin wetness. On the damage side, \citet{davenport2021} used state-level SHELDUS data in a fixed-effects panel and attributed about 36\% of 1988--2017 flood damages to historical precipitation change, mostly from extremely wet conditions. At the county scale, \citet{brody2008} found that rainfall the day before a flood was the strongest physical predictor of property damage in eastern Texas. Damage should therefore be modeled with soil wetness and snow as well as rain, and expected to vary by region.

\paragraph{Exposure and vulnerability.}
How much is in harm's way often matters as much as the storm. \citet{brody2008} found more damage where wetlands had been altered and land paved over, and less where communities had dams or took part in FEMA's Community Rating System. \citet{liu2015} projected Midwest flood damage to rise 13--17\% over about two decades, where county-level damages vary from 2\% to 40\% depending on the location.\citet{tellman2020} showed that rural and socially vulnerable counties suffer more property damage for a given flood intensity. Both works support modeling at the county level for direct insights.

\paragraph{Adaptation.}
Losses also depend on how prepared places are. \citet{kreibich2022} found that risk management generally lowers flood impacts, but not when a new event exceeds what defenses were designed for. \citet{gallagher2013} showed that flood-insurance take-up rises about 9\% after a flood and fades over about nine years, so preparedness follows recent experience. Because wetter places may both flood more and prepare more, adaptation can hide part of the true effect of storms on damage.

\section{Model Description}
\subsection{Problem Setup and Causal Framing}

We follow the standard risk framework of impact forecasting \citep{merz2020}: loss arises when a
\emph{hazard} meets \emph{exposed} property that is \emph{vulnerable} to it. Following the course
framing, we add the county's rain climate and its adaptation, giving the causal structure in
Figure~\ref{fig:dag}.

\begin{table}[H]
\centering
\small
\renewcommand{\arraystretch}{1.2}
\begin{tabularx}{\linewidth}{@{}lL@{}}
\toprule
Symbol & Meaning \\
\midrule
$W$ (storm)        & How much rain fell: the county's wettest 5-day rainfall in the month, in mm. \\
$C$ (rain climate) & How wet the county normally is: the 5-day rainfall it can expect about once every 10 years, in mm. \\
$A$ (adaptation)   & How prepared the county is. Our main measure is flood-insurance policies per 100 homes. \\
$X$ (exposure)     & How much property is in harm's way: the county's built-up area. \\
$Y$ (loss)         & Flood property damage in the county-month combination. \\
\bottomrule
\end{tabularx}
\caption{Variables in the causal diagram (Figure~\ref{fig:dag}).}
\label{tab:dag-vars}
\end{table}

\begin{figure}[H]
\centering
\begin{tikzcd}[column sep=large, row sep=large]
  & C \arrow[dl] \arrow[dr] & \\
  W \arrow[dr] & X \arrow[d] & A \arrow[dl] \\
  & Y &
\end{tikzcd}
\caption{Causal diagram. Rain climate $C$ affects both storm hazard $W$ and
adaptation $A$, opening the backdoor path $W \leftarrow C \rightarrow A \rightarrow Y$.}
\label{fig:dag}
\end{figure}

\subsection{Selecting the Factors}
\label{sec:factors}
Within each part of Figure~\ref{fig:dag}, candidate variables came from the flood-loss literature. A variable was kept only if it met:
\begin{enumerate}
\item \textbf{A physical or economic mechanism}. It must describe a documented way that loss arises (for example, saturated soil turns rain into runoff; floodplain land floods first).
\item \textbf{One variable per mechanism}. Measures of the same mechanism are near-copies that split its effect, so we keep one, and drop one of any pair correlated above 0.8.
\item \textbf{Known before loss}. Yearly measures enter from the previous year, and the climate baseline is built from rainfall only, never from losses.
\item \textbf{Not caused by flood}. Insurance claims, disaster declarations and FEMA disaster payouts are consequences of a flood, so none is used as an input or adjustment variable.
\end{enumerate}

\subsection{Prediction Model}

\paragraph{Hurdle formulation.} Most county-months have no recorded loss, and when loss occurs it
ranges over several orders of magnitude, from a few hundred dollars to hundreds of millions. We
therefore split the prediction into two parts. For county $i$ in month $t$ with inputs
$\mathbf{x}_{it}$, let $L_{it}$ be the loss class: 0 for no damage, and 1--4 for minor (under
\$10k), moderate (\$10k--\$100k), major (\$100k--\$1M) and severe (\$1M or more). Then
\begin{align}
\pi(\mathbf{x}_{it}) &= P(Y_{it} > 0 \mid \mathbf{x}_{it}), \\
q_k(\mathbf{x}_{it}) &= P(L_{it} = k \mid Y_{it} > 0, \mathbf{x}_{it}), \\
P(L_{it} = k \mid \mathbf{x}_{it}) &= \pi(\mathbf{x}_{it}) \, q_k(\mathbf{x}_{it}), \\
P(L_{it} = 0 \mid \mathbf{x}_{it}) &= 1 - \pi(\mathbf{x}_{it}),
\end{align}
for $k = 1, \dots, 4$.
The first part is the probability of any damage; the second is the probability of each loss class if
damage occurs. The first part is fitted on all county-months and the second on damage months
only since a loss class exists only when damage occurs. We predict classes rather than dollar amounts because each class is one order of magnitude,
which matches the precision the data supports. Counties are ranked for emergency resources by the
probability of a severe loss, $\pi(\mathbf{x}_{it}) \, q_4(\mathbf{x}_{it})$. 

\paragraph{Explorer app.} The predictions are presented in a Google Earth Engine app. The
predictions are computed outside Earth Engine and uploaded as a table. A user selects a state (or
all 12) and a month, and the map colors each county by a chosen layer. The prediction layers
show:
\begin{itemize}
  \item the chance of any damage,
  \item rank in the state,
  \item rank in the midwest, 
  \item predicted loss class, 
  \item and the recorded loss class and amount if applicable.
\end{itemize}
The counties in each state's top 5\% are outlined.  Clicking a county lists its chance of damage, its ranks, the predicted loss class with the chance of each class, the recorded loss, and its input values,
with a chart of its predictions over 2017--2023. The predictions use the rain observed in each
month, so they estimate the impact of a known storm rather than forecast one in advance.

\paragraph{Inputs.} Both parts of the hurdle model use the same twelve inputs, one per mechanism (Section~\ref{sec:factors}), plus calendar month for seasonality and state for differences in loss reporting:

\begin{table}[H]
\centering
\small
\rowcolors{2}{gray!12}{white}
\setlength{\aboverulesep}{0pt}
\setlength{\belowrulesep}{0pt}
\renewcommand{\arraystretch}{1.25}
\begin{tabularx}{\linewidth}{@{}lL@{}}
\toprule
Group & Inputs \\
\midrule
Hazard, same month     & monthly rain; wettest 5 days relative to the county's 10-year level; snowmelt; days below freezing \\
Hazard, start of month & root-zone soil wetness (percentile); snowpack \\
Rain climate           & the county's 10-year 5-day rain level \\
Exposure               & built-up surface; impervious surface \\
Susceptibility         & share of land less than 5\,m above drainage; county land area \\
Adaptation             & wetland share \\
Controls               & calendar month; state \\
\bottomrule
\end{tabularx}
\caption{Inputs to the prediction model.}
\label{tab:pred-inputs}
\end{table}

Yearly inputs enter from the previous year. The month's observed rain is used
as an input, so the model assumes the hazard forecast is perfect; in practice
the hazard would come from a forecast.

\paragraph{Main model: LightGBM.} LightGBM builds many small decision trees in sequence, each one
correcting the errors of the trees before it. We chose it because trees capture non-linear
effects and interactions, such as rain falling on wet or frozen ground, without specifying them
in advance. Each part is a gradient-boosted tree ensemble:  a classifier for the probability of any damage, and a four-class classifier for the loss class given damage. We used standard, conservative settings: a learning rate of 0.03, 31 leaves per tree, at least 100 rows per leaf, and 80\%
row and column sampling. Together these keep the model from overfitting the few damage months.
The number of trees was tuned through early stopping (fit on 2000--2012, stop when the error on
2013--2016 no longer improves). The final model is refitted on 2000--2016 with 15\% more trees,
since it sees more data. It is trained on all 12 states, with state as an input.


\paragraph{Comparison model: GLM.} A generalized linear model links the inputs to the outcome
through a single weighted sum, so each input has one coefficient that can be read directly.
We include it because LightGBM's trees have no such coefficients. The GLM shows which inputs
matter and in which direction. Part 1 is a logistic regression for the chance of damage. Part 2
is an ordinal logistic regression, which treats the four loss classes as ordered. Both use a
ridge penalty, chosen on 2013--2016, to keep the coefficients stable.

\paragraph{Baseline.} The baseline uses no inputs: each county's past share of damage months for
that calendar month, and its past share of each loss class. It answers whether the models
learn anything beyond where and when floods have usually happened.

\subsection{Causal Model}
\label{sec:causal}

The prediction model is judged by how accurately it ranks counties. The causal model instead estimates the \emph{storm effect}: how much more likely an extra 25\,mm of rain makes flood damage in a county, and whether that effect is smaller in better-prepared counties. Because wetter counties get bigger storms and
may also prepare more (Figure~\ref{fig:dag}), comparing storms across counties without accounting
for climate and preparedness can misstate the storm effect. We follow the course design of five
nested models, measuring every variable in real world units. In particular, the storm is
rain measured in millimetres rather than relative to the county's normal rainfall levels, since the relative rainfall measure combines
storm and local climate into one number, which hides their separate effects on the outcome.

\paragraph{Data and outcomes.} We use every county-month in 2011--2023, the years with
flood-insurance data, including months with no damage. We model two outcomes: whether the county
recorded \emph{any} flood damage, and whether it recorded a \emph{major} loss of \$100{,}000 or
more.

\paragraph{Variables.} Storm is the effect we estimate; climate may distort or
change it (Table~\ref{tab:adj-vars}). We include all five preparedness measures of the
course design, with insurance as the main one, since it most directly measures flood preparedness.

\begin{table}[H]
\centering
\small
\renewcommand{\arraystretch}{1.15}
\begin{tabularx}{\linewidth}{@{}lL>{\raggedright\arraybackslash}p{2.0cm}@{}}
\toprule
Part of the diagram & What we measure & Effect described per \\
\midrule
$W$ storm        & rain in the month's wettest 5 days (mm)              & 25\,mm more (about 1 inch) \\
$C$ climate      & rain the county expects in its 10-year storm (mm)    & 25\,mm wetter \\
$X$ exposure     & built-up area (km$^2$)                               & twice as much \\
$A$ preparedness & flood-insurance policies per 100 homes (main)        & 1 more per 100 homes \\
                 & median household income (2024 \$)                    & \$10{,}000 more \\
                 & share of homes built before 1980                     & 10 points more \\
                 & share of the county in wetland                       & 5 points  more \\
                 & any community in the Community Rating System         & yes vs no \\
Background       & soil wetness, snowpack and frozen ground at the start of the month;
                   low-lying land; calendar month; state                & -- \\
\bottomrule
\end{tabularx}
\caption{Variables in the causal model. ``Points'' are percentage points: for example, a share of homes built before 1980 rising from 60\% to 70\% is 10 points more.}
\label{tab:adj-vars}
\end{table}


\paragraph{Model.} For each outcome we fit a logistic regression: the chance of damage rises
smoothly with a weighted sum of the variables,
\begin{align}
  P(\text{damage}) &= \frac{1}{1 + e^{-\text{score}}}, \\
  \text{score} &= \beta\,\text{Rain} + \gamma\,\text{Climate} + \lambda\,\text{Built-up} \nonumber \\
               &\quad + \sum_j \alpha_j\,\text{Prep}_j + \text{background} \nonumber \\
               &\quad + \sum_j \delta_j\,(\text{Rain} \times \text{Prep}_j).
\end{align}
Each preparedness measure is centered at its average, so $\beta$ is the storm effect in a county of
average preparedness. The weight $\delta_j$ answers whether the same extra rain does less harm in
counties with more of measure $j$: a negative $\delta_j$ means yes. We read preparedness from
$\delta_j$ rather than $\alpha_j$, because $\alpha_j$ also reflects which counties choose to
prepare; high-risk counties, for example, buy more insurance.

\paragraph{Five models.} We add the parts of Figure~\ref{fig:dag} in stages, on the same
county-months, so that the change in the storm effect shows what each part was hiding
(Table~\ref{tab:m0m4}).

\begin{table}[H]
\centering
\small
\renewcommand{\arraystretch}{1.1}
\begin{tabularx}{\linewidth}{@{}llL@{}}
\toprule
Model & Adds & Compares storms \ldots \\
\midrule
M0 & rain, built-up, background & ignoring climate and preparedness \\
M1 & + climate                  & among similar climates \\
M2 & + preparedness             & among similarly prepared counties \\
M3 & + both                     & among counties similar in both (best estimate) \\
M4 & + rain $\times$ preparedness & and whether it is smaller where better prepared \\
\bottomrule
\end{tabularx}
\caption{The five nested models.}
\label{tab:m0m4}
\end{table}
A change in the storm effect from M0 to M3 shows how much ignoring climate and preparedness
would mislead. Standard errors account for repeated months from the same county.

\paragraph{Interpreting the results.} We report $e^b$: how many times more likely damage becomes when a variable increases by one step. For example, $e^b = 2.4$ for 25\,mm more rain means damage becomes about 2.4 times as likely. $e^b$ multiplies the odds rather than the chance of
damage, but for a rare event like this (about 3\% of county-months) the two are almost the same. 

\begin{comment}
For the main effects we also give the change in chance in percentage points, such as from 2.85\% to about 5.7\%.
\end{comment}

\paragraph{Checks.} We refit with ordinary regression on the yes/no outcomes, whose coefficients
are changes in probability directly; with storm size relative to the county's 10-year level; and with each county's past flood frequency added as a control, since
counties that flooded before buy more insurance.


\section{Experiment}
\subsection{Data}

\paragraph{Study area and outcome.} The 1,055 counties of 12 Midwest states (Illinois, Indiana,
Iowa, Kansas, Michigan, Minnesota, Missouri, Nebraska, North Dakota, Ohio, South Dakota,
Wisconsin), monthly from 2000 to 2023. The outcome is county-month flood property loss from
SHELDUS, in 2024 dollars.

\paragraph{Inputs.} Table~\ref{tab:inputs} lists the inputs, averaged over each county boundary (Census TIGER 2018). Yearly inputs enter from the previous year, so every input except the month's own weather was known before the month began. The 10-year level overlaps the first training years but uses rain only and ends before the test years.

\begin{table}[H]
\centering
\small
\renewcommand{\arraystretch}{1.10}
\begin{tabularx}{\linewidth}{@{}l>{\raggedright\arraybackslash}p{1.5cm}L@{}}
\toprule
Source & Resolution & Variables \\
\midrule
\multicolumn{3}{@{}l}{\textit{Prediction and causal models}} \\
PRISM       & 4 km, daily & rain; days below freezing; wettest 5-day rain and the county's 10-year
                              5-day rain level (fitted on 2000--2009) \\
ERA5-Land   & 11 km, daily and monthly & snowmelt$^\dagger$; snowpack and root-zone soil moisture on day 1 (soil
                              moisture as a percentile vs.\ 1950--1999, same county and month) \\
GHSL        & 100 m, 5-yearly (interpolated) & built-up surface \\
Annual NLCD & 30 m, yearly & paved share$^\dagger$; wetland share \\
MERIT Hydro & 90 m, static & land $<$5 m above drainage \\
\midrule
\multicolumn{3}{@{}l}{\textit{Causal model only}} \\
IPUMS NHGIS & county & median household income (2024 \$); share of homes built before 1980 \\
OpenFEMA    & county, yearly & flood insurance policies per 100 homes; any community in the
                              Community Rating System \\
SHELDUS     & county, monthly & past flood frequency: share of months with damage in the previous
                              10 years (one check only; otherwise outcome only) \\
\bottomrule
\end{tabularx}
\caption{Model inputs by source. The prediction model uses storm size as the wettest 5-day rain
divided by the 10-year level; the causal model uses both in millimetres. Census data from IPUMS
NHGIS \citep{nhgis2026}; flood-insurance data from OpenFEMA \citep{fema2026}, used without FEMA
endorsement. $^\dagger$Prediction model only.}
\label{tab:inputs}
\end{table}


\subsection{Setup}
SHELDUS lists only the months that had a loss. We therefore build a full panel: every county for
every month, 303,840 county-months. The
panel is split by time: train on 2000--2016 and test on 2017--2023. Model settings are chosen by fitting 2000--2012 and scoring 2013--2016.  All 12 states are pooled because a single state has too few damage months to learn the size of the loss.

\subsection{Prediction Results}

\begin{table}[H]
\centering
\small
\renewcommand{\arraystretch}{1.15}
\begin{tabularx}{\linewidth}{@{}Lrrr@{}}
\toprule
Test 2017--2023 & Baseline & GLM & LightGBM \\
\midrule
ROC-AUC (any damage)               & 0.728 & 0.865 & \textbf{0.891} \\
Exact loss class (damage months)   & 36\%  & 38\%  & \textbf{41\%} \\
Within one class                   & 84\%  & \textbf{87\%} & 86\% \\
Top 5\% of counties: share of loss & 24\%  & 61\%  & \textbf{62\%} \\
\bottomrule
\end{tabularx}
\caption{Prediction scores on the test years.}
\label{tab:pred-results}
\end{table}

We found LightGBM to be the more accurate model. To predict the size of the
loss when damage occurs, we first evaluated a regression on log10 dollars. Its typical error was a factor of about 7, which indicated that the data supports order-of-magnitude estimates rather than exact amounts. We therefore model the loss class directly. To compare the two approaches,
we converted each dollar prediction into the class it falls in (for example, a predicted \$50k
loss counts as moderate) and scored both on the same test months. The class model achieves similar
overall accuracy (41\% vs 40\%), but identifies substantially more extreme losses: 34\% of minor
months (vs 16\%) and 33\% of severe months (vs 24\%). Ranking counties by the predicted chance of
a severe loss retains nearly all of the ranking performance of the dollar-based approach (62\% vs
64\% of recorded loss in the top 5\% of counties).

\subsection{Causal Results}
\label{sec:adj-results}

The causal model covers all 164{,}580 county-months in 2011--2023. In an average month, a county
has a 2.85\% chance of recording any flood damage and a 1.05\% chance of a loss of \$100{,}000 or
more. Table~\ref{tab:causal-coefs} gives the coefficients of the full model (M3) and, for
preparedness, of M4. A coefficient $b$ changes the log-odds of damage per step, and $e^b$ is how
many times the odds are multiplied. Because damage is rare, odds and chances are almost the same,
so $e^b = 2.4$ reads as ``about 2.4 times as likely''.




\begin{table*}[t]
\centering
\small
\renewcommand{\arraystretch}{1.15}
\begin{tabularx}{\textwidth}{@{}Lcccc@{}}
\toprule
 & \multicolumn{2}{c}{Any damage} & \multicolumn{2}{c}{Loss of \$100k+} \\
\cmidrule(lr){2-3}\cmidrule(lr){4-5}
Variable (per step) & $b$ (SE) & $e^b$ & $b$ (SE) & $e^b$ \\
\midrule
\multicolumn{5}{@{}l}{\emph{Storm, exposure and climate (M3)}} \\
Rain in wettest 5 days (+25\,mm)       & 0.877 (0.018)* & 2.40 & 0.985 (0.023)* & 2.68 \\
\quad same, without climate or preparedness (M0) & 0.885 (0.017)* & 2.42 & 0.976 (0.022)* & 2.66 \\
Built-up area (doubled)                & 0.230 (0.035)* & 1.26 & 0.366 (0.041)* & 1.44 \\
Climate, same rain (+25\,mm)           & 0.121 (0.056)* & 1.13 & $-0.144$ (0.064)* & 0.87 \\
\midrule
\multicolumn{5}{@{}l}{\emph{Preparedness (M3)}} \\
Flood insurance (+1 per 100 homes)     & 0.129 (0.024)* & 1.14 & 0.065 (0.026)* & 1.07 \\
Median income (+\$10k)                 & $-0.222$ (0.041)* & 0.80 & $-0.149$ (0.040)* & 0.86 \\
Homes built before 1980 (+10 pts)      & $-0.034$ (0.039) & 0.97 & 0.051 (0.042) & 1.05 \\
Wetland share (+5 pts)                 & $-0.089$ (0.034)* & 0.92 & $-0.085$ (0.037)* & 0.92 \\
Any CRS community (yes)                & 0.120 (0.095) & 1.13 & $-0.034$ (0.120) & 0.97 \\
\midrule
\multicolumn{5}{@{}l}{\emph{Does preparedness change the effect of rain? Rain $\times$ measure (M4)}} \\
Rain $\times$ insurance                & $-0.008$ (0.010) & 0.99 & $-0.012$ (0.012) & 0.99 \\
Rain $\times$ income                   & 0.028 (0.012)* & 1.03 & 0.015 (0.014) & 1.02 \\
Rain $\times$ homes before 1980        & 0.008 (0.011) & 1.01 & $-0.017$ (0.013) & 0.98 \\
Rain $\times$ wetland share            & 0.035 (0.012)* & 1.04 & 0.051 (0.016)* & 1.05 \\
Rain $\times$ any CRS community        & 0.027 (0.037) & 1.03 & $-0.051$ (0.041) & 0.95 \\
\midrule
\multicolumn{5}{@{}l}{\emph{Background (M3)}} \\
Soil wetness (+10 percentile points)   & 0.103 (0.007)* & 1.11 & 0.143 (0.011)* & 1.15 \\
Snowpack (doubled)                     & 0.248 (0.021)* & 1.28 & 0.392 (0.028)* & 1.48 \\
Days below freezing (+10)              & $-0.14$ (0.06)* & 0.87 & $-0.04$ (0.08) & 0.96 \\
Low-lying land (+10 pts)               & $-0.034$ (0.018) & 0.97 & 0.014 (0.020) & 1.01 \\
\bottomrule
\end{tabularx}
\caption{Logistic-regression coefficients, 12 states, 2011--2023, with standard errors clustered
by county. $e^b$ is the factor by which one step multiplies the odds (about the chance) of
damage; for the interactions, the factor by which one step of the measure multiplies the effect
of 25\,mm more rain. Models also include calendar month and state. * 95\% interval excludes zero.}
\label{tab:causal-coefs}
\end{table*}

\paragraph{Storm $\rightarrow$ loss: rain is the main driver.} An extra 25\,mm (about an inch) of
rain in a county's wettest five days makes any damage about 2.4 times as likely ($b = 0.877$) and a
\$100k+ loss about 2.7 times as likely ($b = 0.985$). In percentage points, this roughly doubles
the chance of damage, from 2.85\% to about 5.7\%. No other variable comes close: rain is the
dominant cause of damage in a given month, and its effect is larger for big losses than for small
ones.

\paragraph{The backdoor through climate does not distort the rain effect.} Figure~\ref{fig:dag}
warns that comparing storms across counties could mix the storm effect with climate and
preparedness. In practice the rain coefficient hardly moves when they are added: 0.885 without
them (M0) and 0.877 with them (M3) for any damage, and 0.976 and 0.985 for \$100k+ losses. With
rain measured in millimetres, the storm effect is therefore not confounded by climate in our
data. With the course design's storm measure (rain relative to the county's 10-year level), the
storm effect does rise when climate is added; that rise comes from the measure itself, which
combines storm and climate in one number.

\paragraph{Exposure $\rightarrow$ loss: more buildings, more damage.} Doubling a county's
built-up area makes any damage 26\% more likely and a \$100k+ loss 44\% more likely. Exposure
matters more for large losses, as expected: more buildings in the water's path mean more
property to damage once flooding starts.

\paragraph{Climate $\rightarrow$ loss: wet climates have more minor damage but fewer major
losses.} For the same rain, a county whose 10-year storm is 25\,mm larger is 13\% more likely to
have some damage ($b = 0.121$), but 13\% less likely to have a \$100k+ loss ($b = -0.144$). Wet
climates see heavy rain more often, so small damage is more common, but places used to heavy rain
appear to have adapted, so the same storm less often causes a major loss. This is the adaptation
path $C \rightarrow A \rightarrow Y$ in Figure~\ref{fig:dag}, but the effect is small.

\paragraph{Preparedness $\rightarrow$ loss: richer counties and wetlands are damaged less often.}
Each \$10{,}000 of median income makes any damage 20\% less likely and a \$100k+ loss 14\% less
likely, and each 5 points of wetland share makes both 8\% less likely, consistent with money to
protect property and with wetlands storing floodwater. Older homes and Community Rating System
participation show no clear effect. Flood insurance goes the other way: each additional policy per
100 homes makes damage 14\% \emph{more} likely. Insurance does not cause damage; counties at high
flood risk, such as those along major rivers, buy more of it, so this coefficient reflects where
insurance is bought rather than what it does. For this reason we read preparedness from the
interactions below.

\paragraph{Preparedness $\times$ storm: preparation lowers risk but does not weaken rain.} The
interaction terms ask whether the storm effect, about 2.4 times as likely in the average county,
is smaller in better-prepared counties, that is, whether preparation makes heavy rain less
dangerous. It does not. For insurance, rain raises the risk by the same factor at every
level ($e^b = 0.99$). For income and wetland share, the factor is slightly \emph{larger}: each
\$10{,}000 of income multiplies the storm effect by 1.03, and each 5 points more of the county in wetland by 1.04.
What preparation changes is a county's starting chance of damage, not how much heavy rain raises
it: richer counties are damaged less often to begin with, and well-insured counties more often.
Multiplying a lower starting chance by the same factor adds fewer percentage points; for example,
an extra 25\,mm of rain takes a county's chance of damage from 2\% to about 4.7\%, but a county from 4\% to about 9\%.
This is why the same rain adds 2.4\% of chances of damage in richer counties but 3.5\% in poorer ones (90th vs 10th
percentile of income), and 3.1\% points in well-insured counties but 2.7\% in poorly insured ones
(not shown in the table).

\paragraph{Checks.} A linear probability model gives the same conclusions: 25\,mm more rain raises
the chance of any damage by about 4.2 \% (more than the logistic model's 2.9\%, because a straight line fits the curved relationship less closely), and wetter climates see fewer \$100k+ losses. Adding each county's past flood frequency leaves the rain effect unchanged and halves the higher chance of damage in well-insured counties.


\subsection{Limitations}

\begin{itemize}
  \item \textbf{SHELDUS loss data.} In 64\% of SHELDUS damage months the exact dollar value is
        shared with another county, so event totals are often split evenly; 57\% of values are
        round-number estimates; minor losses rose from 19\% of damage months (2000--2012) to 30\%
        (2013--2023); and unreported damage counts as no damage.
  \item \textbf{Upstream river floods.} All inputs describe the county itself, so floods carried
        by rivers from elsewhere are missed. Sarpy County, Nebraska lost \$518M in March 2019 with
        a predicted 5.5\% probability of damage, and a quarter of test damage months were given
        under 5\%.
\item \textbf{Over-prediction.} Fewer damage months were reported after 2016, so the model
        predicts 1.48 times as many as were recorded in 2019--2023. 
  \item \textbf{Observed weather.} The model uses each month's observed rain, so it estimates the
        impact of a known storm. Using only information available before the month, a comparable
        LightGBM's ROC-AUC falls from about 0.90 to 0.76.
  \item \textbf{Associations, not causes.} The causal models control only for measured
        factors, and counties choose how much to adapt. Insurance likely marks of flood-prone  rather than protection, and past flood frequency only partly controls for this; the CRS measure is coarse (only 200
        of 1{,}055 counties take part).
\item \textbf{Reused test years.} The test years were scored several times while we chose the study design, so the test scores are somewhat optimistic.
\end{itemize}


\section{Future Work}
\label{sec:future}

\paragraph{From impact model to forecast.} The model uses each month's observed rain. Feeding it
a hazard forecast, such as precipitation forecasts or the National Water Model \citep{nwm2026}, would complete the
chain from hazard forecast to loss forecast \citep{merz2020} and show how much forecast error adds.

\paragraph{Better data.} Streamflow from USGS gauges \citep{usgs2026} , or rainfall over each county's upstream
catchment, would capture river floods from elsewhere, such as Sarpy County in March 2019. Checking
SHELDUS against NFIP claims \citep{fema2026}, treating losses split across counties as one event, and down-weighting
round-number estimates would reduce noise in the outcome. A Community Rating System measure of
coverage (the share of homes in participating communities), levees from the National Levee
Database \citep{usace2026} and local building codes would improve the adaptation measures.

\paragraph{Stronger causal analysis.} Implementing river-level data, such as distance to major rivers or upstream flow, would separate living in a flood-prone place from buying insurance.

\paragraph{Calibration, scope and uncertainty.} Because reporting shifts over time, both the
probability of damage and the loss classes should be recalibrated each year. Rainfall back to 1951 \citep{durre2022}
would allow training on 1985--2023, roughly doubling the damage months, and applying the model
outside the Midwest would test how well it generalizes. Prediction intervals, for example by
conformal prediction \citep{angelopoulos2021}, would show emergency managers which county rankings are reliable and which
are close calls.

\begin{thebibliography}{}

\bibitem[Angelopoulos and Bates(2021)]{angelopoulos2021}
Angelopoulos, A.N. and Bates, S. (2021)
`A gentle introduction to conformal prediction and distribution-free uncertainty quantification'.
\textit{arXiv preprint} arXiv:2107.07511. \url{https://arxiv.org/abs/2107.07511}

\bibitem[{Apache Software Foundation}(2026)]{arrow2026}
Apache Software Foundation (2026)
\textit{Apache Arrow: PyArrow (Python library)} [software]. \url{https://arrow.apache.org/docs/python/}

\bibitem[Berghuijs et~al.(2016)]{berghuijs2016}
Berghuijs, W.R., Woods, R.A., Hutton, C.J. and Sivapalan, M. (2016)
`Dominant flood generating mechanisms across the United States'.
\textit{Geophysical Research Letters}. 43(9). pp.~4382--4390. \url{https://doi.org/10.1002/2016GL068070}

\bibitem[Brody et~al.(2008)]{brody2008}
Brody, S.D., Zahran, S., Highfield, W.E., Grover, H. and Vedlitz, A. (2008)
`Identifying the impact of the built environment on flood damage in Texas'.
\textit{Disasters}. 32(1). pp.~1--18. \url{https://doi.org/10.1111/j.0361-3666.2007.01024.x}

\bibitem[Davenport et~al.(2021)]{davenport2021}
Davenport, F.V., Burke, M. and Diffenbaugh, N.S. (2021)
`Contribution of historical precipitation change to US flood damages'.
\textit{Proceedings of the National Academy of Sciences}. 118(4). e2017524118. \url{https://doi.org/10.1073/pnas.2017524118}

\bibitem[Durre et~al.(2022)]{durre2022}
Durre, I., Arguez, A., Schreck, C.J., Squires, M.F. and Vose, R.S. (2022)
`Daily high-resolution temperature and precipitation fields for the contiguous United States from
1951 to present'. \textit{Journal of Atmospheric and Oceanic Technology}. 39(12). pp.~1837--1855.
\url{https://doi.org/10.1175/JTECH-D-22-0024.1}

\bibitem[{Federal Emergency Management Agency}(2026)]{fema2026}
Federal Emergency Management Agency (2026)
\textit{OpenFEMA dataset: FIMA NFIP Redacted Policies -- v3}.
Washington, DC: FEMA. Accessed 9 November 2026.
\url{https://www.fema.gov/openfema-data-page/fima-nfip-redacted-policies-v3}

\bibitem[Gallagher(2013)]{gallagher2013}
Gallagher, J. (2013)
\textit{Learning about an infrequent event: Evidence from flood insurance take-up in the US}.
Working paper. Case Western Reserve University, Cleveland, OH. \url{https://www.montana.edu/gallagher/documents/flood_insurance_aeja.pdf}

\bibitem[Google(2026)]{eeapi2026}
Google (2026)
\textit{Earth Engine Python API (earthengine-api)} [software]. \url{https://github.com/google/earthengine-api}

\bibitem[Gorelick et~al.(2017)]{gorelick2017}
Gorelick, N., Hancher, M., Dixon, M., Ilyushchenko, S., Thau, D. and Moore, R. (2017)
`Google Earth Engine: Planetary-scale geospatial analysis for everyone'.
\textit{Remote Sensing of Environment}. 202. pp.~18--27. \url{https://doi.org/10.1016/j.rse.2017.06.031}

\bibitem[Harris et~al.(2020)]{harris2020}
Harris, C.R., Millman, K.J., van der Walt, S.J., Gommers, R., Virtanen, P., Cournapeau, D., et~al. (2020)
`Array programming with NumPy'.
\textit{Nature}. 585. pp.~357--362. \url{https://doi.org/10.1038/s41586-020-2649-2}

\bibitem[Hunter(2007)]{hunter2007}
Hunter, J.D. (2007)
`Matplotlib: A 2D graphics environment'.
\textit{Computing in Science \& Engineering}. 9(3). pp.~90--95. \url{https://doi.org/10.1109/MCSE.2007.55}

\bibitem[Ke et~al.(2017)]{ke2017}
Ke, G., Meng, Q., Finley, T., Wang, T., Chen, W., Ma, W., Ye, Q. and Liu, T.-Y. (2017)
`LightGBM: A highly efficient gradient boosting decision tree'.
\textit{Advances in Neural Information Processing Systems}. 30. pp.~3146--3154.
\url{https://proceedings.neurips.cc/paper/2017/hash/6449f44a102fde848669bdd9eb6b76fa-Abstract.html}

\bibitem[Kreibich et~al.(2022)]{kreibich2022}
Kreibich, H. et~al. (2022)
`The challenge of unprecedented floods and droughts in risk management'.
\textit{Nature}. 608. pp.~80--86. \url{https://doi.org/10.1038/s41586-022-04917-5}

\bibitem[Liu et~al.(2015)]{liu2015}
Liu, J., Hertel, T.W., Diffenbaugh, N.S., Delgado, M.S. and Ashfaq, M. (2015)
`Future property damage from flooding: Sensitivities to economy and climate change'.
\textit{Climatic Change}. 132. pp.~741--749. \url{https://doi.org/10.1007/s10584-015-1478-z}

\bibitem[McKinney(2010)]{mckinney2010}
McKinney, W. (2010)
`Data structures for statistical computing in Python'.
In: \textit{Proceedings of the 9th Python in Science Conference}. pp.~56--61.
\url{https://doi.org/10.25080/Majora-92bf1922-00a}

\bibitem[Merz et~al.(2020)]{merz2020}
Merz, B., Kuhlicke, C., Kunz, M., Pittore, M., Babeyko, A., Bresch, D.N., et~al. (2020)
`Impact forecasting to support emergency management of natural hazards'.
\textit{Reviews of Geophysics}. 58. e2020RG000704. \url{https://doi.org/10.1029/2020RG000704}

\bibitem[{NOAA Office of Water Prediction}(2026)]{nwm2026}
NOAA Office of Water Prediction (2026)
\textit{The National Water Model}. \url{https://water.noaa.gov/about/nwm}

\bibitem[Pedregosa et~al.(2011)]{pedregosa2011}
Pedregosa, F., Varoquaux, G., Gramfort, A., Michel, V., Thirion, B., Grisel, O., et~al. (2011)
`Scikit-learn: Machine learning in Python'.
\textit{Journal of Machine Learning Research}. 12. pp.~2825--2830. \url{https://jmlr.org/papers/v12/pedregosa11a.html}

\bibitem[{Python Software Foundation}(2026)]{python2026}
Python Software Foundation (2026)
\textit{Python Language Reference} [software]. \url{https://www.python.org/}

\bibitem[Schroeder et~al.(2026)]{nhgis2026}
Schroeder, J., Van Riper, D., Manson, S., Cooper, G., Krause, Z., Kugler, T., Zhu, T. and
Ruggles, S. (2026)
\textit{IPUMS National Historical Geographic Information System: Version 21.0} [dataset].
Minneapolis, MN: IPUMS. \url{https://doi.org/10.18128/D050.V21.0}

\bibitem[Seabold and Perktold(2010)]{seabold2010}
Seabold, S. and Perktold, J. (2010)
`Statsmodels: Econometric and statistical modeling with Python'.
In: \textit{Proceedings of the 9th Python in Science Conference}. pp.~92--96.
\url{https://doi.org/10.25080/Majora-92bf1922-011}

\bibitem[Slater and Villarini(2016)]{slater2016}
Slater, L.J. and Villarini, G. (2016)
`Recent trends in U.S. flood risk'.
\textit{Geophysical Research Letters}. 43(24). pp.~12428--12436. \url{https://doi.org/10.1002/2016GL071199}

\bibitem[Smith et~al.(2026)]{patsy2026}
Smith, N.J., et~al. (2026)
\textit{Patsy: Describing statistical models in Python} [software]. \url{https://patsy.readthedocs.io/}

\bibitem[Tellman et~al.(2020)]{tellman2020}
Tellman, B., Schank, C., Schwarz, B., Howe, P.D. and de Sherbinin, A. (2020)
`Using disaster outcomes to validate components of social vulnerability to floods: Flood deaths and property damage across the USA'.
\textit{Sustainability}. 12(15). 6006. \url{https://doi.org/10.3390/su12156006}

\bibitem[{U.S. Army Corps of Engineers}(2026)]{usace2026}
U.S. Army Corps of Engineers (2026)
\textit{National Levee Database}. \url{https://levees.sec.usace.army.mil/}

\bibitem[{U.S. Geological Survey}(2026)]{usgs2026}
U.S. Geological Survey (2026)
\textit{National Water Information System}. \url{https://waterdata.usgs.gov/nwis}

\bibitem[Virtanen et~al.(2020)]{virtanen2020}
Virtanen, P., Gommers, R., Oliphant, T.E., Haberland, M., Reddy, T., Cournapeau, D., et~al. (2020)
`SciPy 1.0: Fundamental algorithms for scientific computing in Python'.
\textit{Nature Methods}. 17. pp.~261--272. \url{https://doi.org/10.1038/s41592-019-0686-2}

\end{thebibliography}
\begin{itemize}
    \item Code development and report revisions were assisted by Claude; however, all references and links were independently reviewed and verified for accuracy.
\end{itemize}


\section{Links}

\subsection*{Github Repository}
\begin{itemize}
    \item \url{https://github.com/AAA-zzZZ/CSE472-Project-2.git}
\end{itemize}

\subsection*{Google Earth Engine App}
\begin{itemize}
    \item \textbf{App: }
    \url{ https://code.earthengine.google.com/?accept_repo=users/ezheng3/midwest_flood_impact}
    \item \textbf{Assets for App: }\url{https://code.earthengine.google.com/?asset=projects/cse472-project-2/assets/midwest_flood_app_table}
    \begin{itemize}
    \item Access for the assets have been granted for Ali Beigi and Saketh Vishnubhatla
    \item Outside Access has to be granted manually
    \end{itemize}
\end{itemize}
\subsection*{Google Earth Engine Resources}

\begin{itemize}
    \item \textbf{PRISM (Climate Data): } \url{https://developers.google.com/earth-engine/datasets/catalog/OREGONSTATE_PRISM_ANd#bands}
    \item \textbf{ERA5-Land: }
    \url{https://developers.google.com/earth-engine/datasets/catalog/ECMWF_ERA5_LAND_MONTHLY_AGGR}\textbf{ and  }\url{https://developers.google.com/earth-engine/datasets/catalog/ECMWF_ERA5_LAND_DAILY_AGGR}
    \item \textbf{GHSL (Global Human Settlement Layer): } \url{https://developers.google.com/earth-engine/datasets/catalog/JRC_GHSL_P2023A_GHS_BUILT_S}
    \item \textbf{Annual NLCD: }\url{https://gee-community-catalog.org/projects/annual_nlcd/}
    \item \textbf{MERIT Hydro: }\url{https://developers.google.com/earth-engine/datasets/catalog/MERIT_Hydro_v1_0_1}
\end{itemize}


\end{document}